\chapter{Guided Reading: Better Than Square Root}
\label{ch:chaos}

\begin{center}\itshape
Flip a fair coin for every prime, multiply,\\
and add up what you get.\\
The sum is smaller than it has any right to be.
\end{center}

Here is an experiment that anyone with a computer can run. For each prime $p$, choose a random point $f(p)$ on the unit circle, independently. Extend $f$ to all positive integers by multiplicativity,
$f(n)=\prod f(p)^{a_p}$ when $n=\prod p^{a_p}$, and add up: $S(x)=\sum_{n\le x}f(n)$. This is a \emph{Steinhaus random multiplicative function}, the model of Chapter~\ref{ch:oc-multiplicative}. Its values are
orthogonal, $\mathbb E\,f(n)\overline{f(m)}=\mathbf 1_{n=m}$, so $\mathbb E|S(x)|^2=\lfloor x\rfloor$, and one would bet that $S(x)$ is typically of size $\sqrt x$, like a sum of $x$ independent
unit vectors. The bet is wrong. In 2010 Henry Helson conjectured that $\mathbb E|S(x)|=o(\sqrt x)$, and in 2020 Adam Harper proved it, with the precise order
\[
\mathbb E|S(x)|\asymp\frac{\sqrt x}{(\log\log x)^{1/4}}.
\]
The correction is tiny: $(\log\log x)^{1/4}$ is less than $2$ for every $x$ you will ever compute with. But it is there, and the reason it is there is one of the most beautiful ideas of recent probability: \emph{critical multiplicative chaos}.

Harper's paper, \emph{Moments of random multiplicative functions, I: low moments, better than squareroot cancellation, and critical multiplicative chaos} (Forum of Mathematics, Pi, 2020), is eighty pages long. This chapter shows how to read it, maps its ingredients to the places in this series where they are proved, and proves the heart of the matter in a model where everything can be done in a few pages: a branching random walk.

\section{How to read an eighty-page paper}

The method of Chapter~\ref{ch:reading} still applies, but at this length the first pass needs a plan.
\begin{enumerate}
\item \emph{The introduction is a map.} Harper's introduction states the main theorems: for $0\le q\le1$, uniformly,
\[
\mathbb E|S(x)|^{2q}\asymp\Bigl(\frac x{1+(1-q)\sqrt{\log\log x}}\Bigr)^q,
\]
in the Steinhaus case, with an analogue for the Rademacher model ($f(p)=\pm1$, supported on squarefree $n$). It then explains, in words, why the answer should be what it is, and where the difficulty lies. Read it twice.
\item \emph{Find the heuristic section and read it before any proof.} The paper explains the connection with multiplicative chaos and with the branching random walk informally before doing anything rigorous; that is the section that tells you what the eighty pages are for.
\item \emph{Take the technical lemmas on trust on a first reading.} Many lemmas are estimates for sums over primes, Euler products and Gaussian approximations. Read their statements, check that they are plausible (often by a two-line heuristic), and move on.
\item \emph{Identify the two or three steps that are new.} Here they are: the reduction of $\mathbb E|S(x)|$ to an integral of an Euler product, the conditioning on the small primes, and a ballot-type estimate for a random walk that keeps the Euler product below a barrier. The rest is careful, hard, but standard analytic number theory.
\end{enumerate}

\section{The map to the series}

Each ingredient has a home earlier in the series.
\begin{itemize}
\item \emph{The models and their first moments.} Orthogonality, the almost sure bound of Wintner, and the fourth moment $\mathbb E|S(x)|^4\sim\frac{12}{\pi^2}x^2\log x$ are stated in Chapter~\ref{ch:oc-multiplicative}. The fourth moment is larger than the Gaussian value $2x^2$: high moments are dominated by rare large values. Harper's theorem says that the low moments are dominated by the opposite, by the absence of large values.
\item \emph{From partial sums to an Euler product.} Plancherel's theorem for Dirichlet series (Volume~14, Lemma~9.4) turns the mean square of partial sums into an integral of a Dirichlet series on a vertical line. Applied conditionally, it shows that $|S(x)|^2$ is, on average over the large primes, comparable to $\frac x{\log x}\int|F(\frac12+\frac1{\log x}+it)|^2\frac{\dd t}{|\frac12+it|^2}$, where $F(s)=\prod_{p\le x}(1-f(p)p^{-s})^{-1}$ is the random Euler product over the small primes.
\item \emph{Conditioning on the small primes.} Given the values of $f$ on primes up to $\sqrt x$, the contribution of the large primes is approximately Gaussian with a random variance $V(x)$: this conditional approximation is stated in Chapter~\ref{ch:oc-multiplicative}. Hence $\mathbb E|S(x)|\approx\mathbb E\sqrt{V(x)}$, a \emph{half moment} of a random variance.
\item \emph{The random variance is a chaos.} By the previous item, $V(x)$ is essentially the integral of $|F|^2$ over $t$. The logarithm $\log|F(\frac12+it)|=\sum_p\Re\frac{f(p)}{p^{1/2+it}}+O(1)$ is a sum of independent pieces, and two points $t,t'$ see the same primes $p<e^{1/|t-t'|}$ and different ones above: $\log|F|$ is a \emph{log-correlated field}, exactly like the logarithm of the characteristic polynomial of a random unitary matrix (Volume~12, Proposition~8.3, and the multiplicative chaos of Volume~12, Theorem~8.4). The integral of $|F|^2$ is then a multiplicative chaos at the \emph{critical} parameter, and its half moment is what we need.
\end{itemize}

So the whole question becomes: why is the half moment of a critical chaos small? The cleanest place to see the answer is a tree.

\section{The model: a critical branching random walk}

Think of the scales of primes as generations. Group the primes into blocks $(e^{e^{k-1}},e^{e^k}]$; each block contributes to $\log|F(\frac12+it)|$ a term of variance about $\frac12$, by Mertens's theorem. A point $t$ at resolution $e^{-k}$ is determined by its first $k$ blocks, and two points at distance about $e^{-j}$ share the first $j$ blocks and then go their separate ways: that is a tree of depth $n\approx\log\log x$, with independent increments along its edges. The following model keeps exactly this structure and nothing else.

\begin{definition}\label{def:brw}
In the \emph{critical branching random walk}, every individual, independently, has a random number $N^+\sim\mathrm{Poisson}(\frac1{2e})$ of children placed one step up and $N^-\sim\mathrm{Poisson}(\frac e2)$ of children placed one step down. The ancestor sits at $0$. For an individual $v$ of generation $|v|=n$ write $V(v)$ for its position, $v_k$ for its ancestor in generation $k$, and put
\[
W_n=\sum_{|v|=n}e^{V(v)}.
\]
\end{definition}

The numbers $\frac1{2e}$ and $\frac e2$ are chosen so that the two conditions of criticality hold: the expected sum of $e^{\text{step}}$ over the children of one individual is $\frac1{2e}\,e+\frac e2\,e^{-1}=1$, and the expected sum of $\text{step}\cdot e^{\text{step}}$ is $\frac12-\frac12=0$. The first makes $W_n$ a martingale with $\mathbb EW_n=1$; the second puts it at the critical point. (The mean number of children is $\frac1{2e}+\frac e2\approx1.54$, so the population grows exponentially when it survives.) The quantity $W_n$ plays the role of $V(x)/x$, and $n$ the role of $\log\log x$.

\begin{lemma}[many-to-one]\label{lem:manytoone}
For every $n\ge1$ and every function $G:\Z^n\to[0,\infty)$,
\[
\mathbb E\sum_{|v|=n}e^{V(v)}G\bigl(V(v_1),\dots,V(v_n)\bigr)=\mathbb E\,G(S_1,\dots,S_n),
\]
where $(S_k)$ is the simple random walk started at $0$.
\end{lemma}

\begin{proof}
Induction on $n$. Condition on the first generation. Given its individuals and their positions $V(c)=\pm1$, the subtrees rooted at them are independent copies of the whole process. The descendants of $c$ in generation $n$ are the individuals $w$ of generation $n-1$ of its subtree, with $V(v)=V(c)+V'(w)$, so $e^{V(v)}=e^{V(c)}e^{V'(w)}$. Let
\[
H(y)=\mathbb E\,G(y,y+S_1,\dots,y+S_{n-1})
\]
(for $n=1$, $H(y)=G(y)$). By the induction hypothesis applied to the subtree of $c$, the conditional expectation of the contribution of the descendants of $c$ is $e^{V(c)}H(V(c))$. It remains to compute $\mathbb E\sum_{|c|=1}e^{V(c)}H(V(c))$. A Poisson number of children with mean $m$ contributes $m$ times the value in expectation, so this is
\[
\tfrac1{2e}\,e\,H(1)+\tfrac e2\,e^{-1}H(-1)=\tfrac12H(1)+\tfrac12H(-1)=\mathbb E\,H(S_1),
\]
which is $\mathbb E\,G(S_1,\dots,S_n)$ by the Markov property of the walk.
\end{proof}

With $G=1$ the lemma gives $\mathbb EW_n=1$. Two more facts about the simple random walk are needed, both consequences of the reflection principle of Volume~4.

\begin{lemma}[ballot bound]\label{lem:ballotbound}
For integers $L\ge0$ and $n\ge1$, $P(S_k\le L\text{ for all }k\le n)\le\dfrac{2L+2}{\sqrt n}$.
\end{lemma}

\begin{proof}
By Volume~4, Proposition~7.5, $P(\max_{k\le n}S_k\ge r)=P(S_n\ge r)+P(S_n>r)$ for $r\ge1$. With $r=L+1$ and the symmetry of $S_n$,
\[
P\bigl(\max_{k\le n}S_k\le L\bigr)=1-P(S_n\ge L+1)-P(S_n\ge L+2)=P(-L-1\le S_n\le L),
\]
which is a sum of $2L+2$ point probabilities. Each is at most $a_n=\binom n{\lfloor n/2\rfloor}2^{-n}$. We show $a_n\le n^{-1/2}$. Let $b_m=\binom{2m}m4^{-m}$; then $b_{m+1}=b_m\frac{2m+1}{2m+2}$, and $b_m^2\le\frac1{3m+1}$ by induction, since $b_0=1$ and
\[
\Bigl(\frac{2m+1}{2m+2}\Bigr)^2\le\frac{3m+1}{3m+4}\iff12m^3+28m^2+19m+4\le12m^3+28m^2+20m+4.
\]
For $n=2m$, $a_n=b_m\le(3m+1)^{-1/2}\le n^{-1/2}$. For $n=2m+1$, $a_n=\frac12\binom{2m+1}m4^{-m}=b_m\frac{2m+1}{2m+2}\le b_m\le(3m+1)^{-1/2}\le n^{-1/2}$.
\end{proof}

\begin{lemma}[no high individuals]\label{lem:nohigh}
Let $M_n=\max\{V(v):|v|\le n\}$. For every integer $L\ge0$, $P(M_n\ge L+1)\le e^{-(L+1)}$.
\end{lemma}

\begin{proof}
If $M_n\ge L+1$, some individual $v$ with $|v|=k\le n$ is the \emph{first} on its line of descent to reach $L+1$: $V(v)=L+1$ (the walk moves by $\pm1$) and $V(v_j)\le L$ for $j<k$. Let $N$ be the number of such individuals. By Lemma~\ref{lem:manytoone} with $G=e^{-y_k}\mathbf 1\{y_k=L+1,\ y_j\le L\ (j<k)\}$,
\[
\mathbb EN=\sum_{k=1}^n e^{-(L+1)}P\bigl(S_k=L+1,\ S_j\le L\ (j<k)\bigr)\le e^{-(L+1)},
\]
since the events for different $k$ are disjoint. So $P(M_n\ge L+1)\le P(N\ge1)\le\mathbb EN$.
\end{proof}

\begin{theorem}[better than square root, in the tree]\label{thm:halfmoment}
For every $n\ge1$, $\mathbb E\bigl[W_n^{1/2}\bigr]\le\dfrac{C}{n^{1/4}}$, with $C=\sum_{L\ge0}\sqrt{2L+2}\,e^{-L/2}<5.4$. In particular $W_n\to0$ in probability, although $\mathbb EW_n=1$ for every $n$.
\end{theorem}

\begin{proof}
Since the ancestor sits at $0$, $M_n$ is an integer $\ge0$, and
\[
\mathbb E\bigl[W_n^{1/2}\bigr]=\sum_{L\ge0}\mathbb E\bigl[W_n^{1/2}\mathbf 1\{M_n=L\}\bigr].
\]
On the event $\{M_n=L\}$, every line of descent stays at or below $L$, so $W_n=W_n^{(L)}$, where $W_n^{(L)}=\sum_{|v|=n}e^{V(v)}\mathbf 1\{V(v_k)\le L\text{ for all }k\le n\}$. By the Cauchy--Schwarz inequality,
\[
\mathbb E\bigl[W_n^{1/2}\mathbf 1\{M_n=L\}\bigr]\le\bigl(\mathbb EW_n^{(L)}\bigr)^{1/2}P(M_n=L)^{1/2}.
\]
By Lemma~\ref{lem:manytoone} and Lemma~\ref{lem:ballotbound}, $\mathbb EW_n^{(L)}=P(S_k\le L\ \forall k\le n)\le\frac{2L+2}{\sqrt n}$. By Lemma~\ref{lem:nohigh}, $P(M_n=L)\le P(M_n\ge L)\le e^{-L}$ (trivially for $L=0$). Summing,
\[
\mathbb E\bigl[W_n^{1/2}\bigr]\le\sum_{L\ge0}\Bigl(\frac{2L+2}{\sqrt n}\Bigr)^{1/2}e^{-L/2}=\frac C{n^{1/4}}.
\]
The series converges, and its value is below $5.4$ (a direct computation gives $5.39\ldots$). Finally, by Markov's inequality, $P(W_n>\varepsilon)\le\varepsilon^{-1/2}\mathbb E[W_n^{1/2}]\to0$.
\end{proof}

Look at what the proof says. The mean $\mathbb EW_n=1$ is carried by rare trees in which some line of descent climbs high: there $W_n$ is huge. On the typical tree, all lines stay below a barrier at height $O(1)$, and the ballot theorem says that only a proportion $n^{-1/2}$ of the mass survives such a constraint. The half moment sees the typical tree, so it sees $n^{-1/4}$. Figure~\ref{fig:barrier} shows a sample.

\begin{figure}[ht]\centering\fig{ch06-barrier}
\caption{A critical branching random walk over $11$ generations. The highest individual ever reached height $M=2$ (dashed); the four lineages carrying most of $W_{11}=\sum e^{V(v)}$ are highlighted. Lineages that climb above the barrier are rare, and in the typical tree the mass that survives below it is a proportion $n^{-1/2}$ of the mean.}\label{fig:barrier}
\end{figure}

\begin{remark}[the lower bound, and the truth]\label{rem:brwlower}
The bound of Theorem~\ref{thm:halfmoment} is sharp: $\mathbb E[W_n^{1/2}]\ge cn^{-1/4}$ for some $c>0$. The proof applies the inequality $\mathbb EZ^{1/2}\ge(\mathbb EZ)^{3/2}/(\mathbb EZ^2)^{1/2}$ (Exercise~\ref{ex:holderlower}) to a version of $W_n^{(0)}$ restricted to lines that end not too far below the barrier, and bounds $\mathbb EZ^2$ by a \emph{many-to-two} lemma and a ballot theorem for random walk bridges; it is longer but uses nothing new. Much more is true: $\sqrt n\,W_n$ converges in probability to a positive multiple of the limit of the \emph{derivative martingale} on the event of survival (Aïdékon and Shi, 2014, after Hu and Shi, 2009).
\end{remark}

\section{Back to the integers}

The dictionary between the tree and the random multiplicative function reads as follows.
\begin{itemize}
\item generation $k$ $\leftrightarrow$ the block of primes $(e^{e^{k-1}},e^{e^k}]$, so $n\leftrightarrow\log\log x$;
\item the position $V(v)$ $\leftrightarrow$ $2\log|F|$ at a point $t$, computed with the primes up to $e^{e^k}$, minus its mean;
\item $W_n$ $\leftrightarrow$ the normalised integral $\frac{\log x}x\,V(x)\approx\int|F(\frac12+\frac1{\log x}+it)|^2\frac{\dd t}{|\frac12+it|^2}$, with $\mathbb E=1$ up to constants;
\item Lemma~\ref{lem:manytoone} $\leftrightarrow$ the Girsanov (tilting) step: weighting by $|F|^2$ turns the sum over primes into a random walk with zero drift;
\item Theorem~\ref{thm:halfmoment} $\leftrightarrow$ Harper's estimate
\[
\mathbb E|S(x)|\approx\sqrt x\,\mathbb E\Bigl[\Bigl(\frac{\log x}xV(x)\Bigr)^{1/2}\Bigr]\asymp\frac{\sqrt x}{(\log\log x)^{1/4}}.
\]
\end{itemize}

\begin{remark}[Harper's theorem]\label{rem:harper2020}
Harper proves, for Steinhaus and Rademacher random multiplicative functions, that $\mathbb E|S(x)|\asymp\sqrt x/(\log\log x)^{1/4}$, and more generally the formula for the $2q$-th moments stated above. The architecture of the proof is the one of this chapter: (1) condition on the small primes, reduce to the half moment of the random variance (the conditional central limit theorem stated in Chapter~\ref{ch:oc-multiplicative}); (2) express the variance as an integral of the Euler product, by Plancherel; (3) prove the upper bound by restricting to the event that the partial Euler products stay below a barrier at every scale, which costs the ballot factor $(\log\log x)^{-1/2}$ in expectation, as in Theorem~\ref{thm:halfmoment}; (4) prove the lower bound by a second moment computation on a barrier event. The difficulty, absent in the tree, is that the Euler product is a continuous field in $t$ and the increments over prime blocks are only approximately Gaussian and independent; most of the eighty pages are spent making the tree picture honest.
\end{remark}

The story continued quickly. A first corollary: $S(x)/\sqrt x\to0$ in probability, so the full sums have no central limit theorem, in contrast with the conditional approximation of Chapter~\ref{ch:oc-multiplicative}. And Harper has since applied the same ideas to Dirichlet character sums and to sums of $n^{it}$, showing that their typical size is also $o(\sqrt x)$. That is the mark of a frontier paper: its main idea is portable.

\begin{think}
In Theorem~\ref{thm:halfmoment} the exponent $\frac14$ is half of the ballot exponent $\frac12$. What exponent would you get for $\mathbb E[W_n^q]$ with $0<q<1$? (Exercise~\ref{ex:qmoment}.) Why does the argument say nothing for $q>1$?
\end{think}

\begin{beauty}
A sum of numbers on the unit circle, one for each integer, is secretly a tree of random walks, and its typical size is decided by whether a walk can stay below a line. The ballot problem of Bertrand, posed in 1887 about counting votes, decides the size of random sums over the integers. Mathematics rarely shows its unity so plainly.
\end{beauty}

\section{Exercises}

\begin{exercise}[$\star$]\label{ex:criticalparams}
In Definition~\ref{def:brw}, replace the means by $a$ (up) and $b$ (down). Show that $\mathbb EW_n=1$ for all $n$ and the many-to-one walk has zero drift if and only if $a=\frac1{2e}$ and $b=\frac e2$.
\end{exercise}

\begin{exercise}[$\star$]\label{ex:martingale}
Show that $(W_n)$ is a martingale for the filtration of the first $n$ generations.
\end{exercise}

\begin{exercise}[$\star\star$]\label{ex:qmoment}
Adapt the proof of Theorem~\ref{thm:halfmoment} to show that $\mathbb E[W_n^q]\le C_qn^{-q/2}$ for $0<q<1$. (Use Hölder's inequality in place of Cauchy--Schwarz.)
\end{exercise}

\begin{exercise}[$\star\star$]\label{ex:holderlower}
Show that $\mathbb EZ^{1/2}\ge(\mathbb EZ)^{3/2}/(\mathbb EZ^2)^{1/2}$ for every random variable $Z\ge0$ with $0<\mathbb EZ^2<\infty$.
\end{exercise}

\begin{exercise}[$\star\star$]\label{ex:subcritical}
For $\beta>0$, give each individual $\mathrm{Poisson}(\frac12e^{-\beta})$ children one step up and $\mathrm{Poisson}(\frac12e^{\beta})$ one step down. Show that $\sum_{|v|=n}e^{\beta V(v)}$ has mean $1$ and that its many-to-one walk is the simple random walk. For $\theta>0$, show that $\mathbb E\sum_{|v|=n}e^{\theta V(v)}=\cosh(\theta-\beta)^n$, and that along the walk of the weights $e^{\theta V}$ the normalised log-weight $\theta V-n\log\cosh(\theta-\beta)$ has drift $\theta\tanh(\theta-\beta)-\log\cosh(\theta-\beta)$. Show that this drift vanishes at $\theta=\beta$ and is negative for $0<\theta<\beta$.
\end{exercise}

\begin{exercise}[$\star$]\label{ex:fourthmomentgap}
Compute $\mathbb EW_n^2$ for small $n$ (by a numerical calculation, or exactly for $n=1$) and show that it tends to infinity. Compare with Theorem~\ref{thm:halfmoment}: the second moment and the half moment of $W_n$ go in opposite directions, as the fourth and the first moments of $S(x)/\sqrt x$ do.
\end{exercise}

\section{Hints}

\begin{hint}{ex:criticalparams}
Compute $\mathbb E\sum_{|c|=1}e^{V(c)}$ and $\mathbb E\sum_{|c|=1}V(c)e^{V(c)}$.
\end{hint}

\begin{hint}{ex:martingale}
Condition on generation $n$ and use the first computation of the previous hint for each individual.
\end{hint}

\begin{hint}{ex:qmoment}
Hölder's inequality gives
\[
\mathbb E[X\mathbf 1_A]\le(\mathbb EX^{1/q})^q P(A)^{1-q}
\]
for $X\ge0$; take $X=(W_n^{(L)})^q$.
\end{hint}

\begin{hint}{ex:holderlower}
Write $Z=Z^{1/3}Z^{2/3}$ and apply Hölder with exponents $\frac32$ and $3$.
\end{hint}

\begin{hint}{ex:subcritical}
Repeat the proof of Lemma~\ref{lem:manytoone} with general weights; the step of the walk for the weights $e^{\theta V}$ is $+1$ with probability proportional to $\frac12e^{-\beta}e^{\theta}$.
\end{hint}

\begin{hint}{ex:fourthmomentgap}
For $n=1$, $W_1=N^+e+N^-e^{-1}$ with independent Poisson variables.
\end{hint}

\section{Solutions}

\begin{solution}{ex:criticalparams}
\idea Two equations, two unknowns.
\writeup For one individual,
\[
\mathbb E\sum_{|c|=1}e^{V(c)}=ae+be^{-1},\qquad\mathbb E\sum_{|c|=1}V(c)e^{V(c)}=ae-be^{-1}.
\]
The first must be $1$ (then $\mathbb EW_n=1$ for all $n$, by the branching property) and the second $0$ (it is the drift of the many-to-one walk, whose step is $+1$ with probability $ae$ and $-1$ with probability $be^{-1}$). So $ae=be^{-1}=\frac12$.
\end{solution}

\begin{solution}{ex:martingale}
\idea Each individual's children reproduce its own weight on average.
\writeup Given the first $n$ generations,
\[
W_{n+1}=\sum_{|v|=n}e^{V(v)}\sum_{c\text{ child of }v}e^{V(c)-V(v)},
\]
and the inner sums are independent of the past with mean $\frac1{2e}e+\frac e2e^{-1}=1$. So $\mathbb E[W_{n+1}\mid\mathcal F_n]=W_n$. ($W_n$ is integrable since $\mathbb EW_n=1$.)
\end{solution}

\begin{solution}{ex:qmoment}
\idea Same decomposition, with Hölder.
\writeup On $\{M_n=L\}$, $W_n=W_n^{(L)}$. By Hölder with exponents $\frac1q$ and $\frac1{1-q}$, $\mathbb E[(W_n^{(L)})^q\mathbf 1\{M_n=L\}]\le(\mathbb EW_n^{(L)})^qP(M_n=L)^{1-q}\le\bigl(\frac{2L+2}{\sqrt n}\bigr)^qe^{-(1-q)L}$. Summing over $L$, $\mathbb E[W_n^q]\le C_qn^{-q/2}$ with $C_q=\sum_L(2L+2)^qe^{-(1-q)L}<\infty$. As $q\to1$, $C_q\to\infty$, consistent with $\mathbb EW_n=1$.
\end{solution}

\begin{solution}{ex:holderlower}
\idea Interpolate the first moment between the half and the second moments.
\writeup By Hölder's inequality with exponents $\frac32$ and $3$,
\[
\mathbb EZ=\mathbb E[Z^{1/3}Z^{2/3}]\le(\mathbb EZ^{1/2})^{2/3}(\mathbb EZ^2)^{1/3},
\]
since $(Z^{1/3})^{3/2}=Z^{1/2}$ and $(Z^{2/3})^3=Z^2$. Raise to the power $\frac32$ and rearrange.
\end{solution}

\begin{solution}{ex:subcritical}
\idea Criticality is the point where the normalised log-weight along the tilted walk has no drift.
\writeup For one individual, $\mathbb E\sum_ce^{\theta V(c)}=\frac12e^{-\beta}e^\theta+\frac12e^\beta e^{-\theta}=\cosh(\theta-\beta)$, which is $1$ for $\theta=\beta$; by the branching property the mean over generation $n$ is $\cosh(\theta-\beta)^n$. The proof of Lemma~\ref{lem:manytoone} goes through with $e^{V}$ replaced by $e^{\theta V}/\cosh(\theta-\beta)$ per generation, and the walk steps $+1$ with probability $\frac{e^{\theta-\beta}}{2\cosh(\theta-\beta)}$; for $\theta=\beta$ this is $\frac12$. The mean step is $\tanh(\theta-\beta)$, so the log-weight $\theta V-n\log\cosh(\theta-\beta)$ has drift $\theta\tanh(\theta-\beta)-\log\cosh(\theta-\beta)$, zero at $\theta=\beta$; for $0<\theta<\beta$ both terms are negative. Biggins's theorem (which we do not prove) says that the normalised martingale has a positive limit with mean $1$ exactly when this drift is negative, the \emph{subcritical} regime; the model of the chapter is the boundary case $\theta=\beta=1$, and there the limit is $0$ (Theorem~\ref{thm:halfmoment}).
\end{solution}

\begin{solution}{ex:fourthmomentgap}
\idea High moments are dominated by rare high individuals.
\writeup For $n=1$, $\mathbb EW_1^2=\operatorname{Var}W_1+1=\frac1{2e}e^2+\frac e2e^{-2}+1=\frac e2+\frac1{2e}+1=1+\cosh1\approx2.54$. By the branching property, $\mathbb EW_{n+1}^2=\mathbb E W_n^2+\mathbb E\sum_{|v|=n}e^{2V(v)}\operatorname{Var}W_1$, and $\mathbb E\sum_{|v|=n}e^{2V(v)}=\mathbb Ee^{S_n}=(\cosh1)^n$ by many-to-one, so $\mathbb EW_n^2=1+\cosh1\cdot\frac{(\cosh1)^n-1}{\cosh1-1}$ grows exponentially. The half moment tends to $0$, the second moment to infinity: the same split as between $\mathbb E|S(x)|\ll\sqrt x$ and $\mathbb E|S(x)|^4\gg x^2$.
\end{solution}

\section{Sources}

\begin{itemize}
\item A.~J.~Harper, Moments of random multiplicative functions, I: low moments, better than squareroot cancellation, and critical multiplicative chaos, \emph{Forum Math.\ Pi} 8 (2020), e1; arXiv:1703.06654. The paper of this chapter. Read the introduction and the heuristic discussion first.
\item H.~Helson, Hankel forms, \emph{Studia Math.} 198 (2010), 79--84: the conjecture $\mathbb E|S(x)|=o(\sqrt x)$.
\item A.~J.~Harper, Moments of random multiplicative functions, II: high moments, \emph{Algebra Number Theory} 13 (2019): the counterpart for $q>1$.
\item Y.~Hu, Z.~Shi, Minimal position and critical martingale convergence in branching random walks, and directed polymers on disordered trees, \emph{Ann.\ Probab.} 37 (2009); E.~Aïdékon, Z.~Shi, The Seneta--Heyde scaling for the branching random walk, \emph{Ann.\ Probab.} 42 (2014).
\item Z.~Shi, \emph{Branching Random Walks}, \'Ecole d'\'Et\'e de Probabilit\'es de Saint-Flour XLII, Lecture Notes in Math.\ 2151, Springer, 2015: the many-to-one lemma and the barrier method.
\end{itemize}
